

Many sources of data have natural relational structure which is well modelled by graphs, such as social networks, urban networks, protein-protein interaction networks and blood flow in the brain. Inevitably, some of this `graph data' will be noisy and incomplete. The field of Graph Signal Processing extends the powerful tools of signal processing to the realm of graph data, and to tackle missing data has extended sampling and reconstruction tools. Many studies have been written on designing sampling schemes, however few give guidance on sample size -- how much data to collect -- or how the task at hand affects the handling of missing data. This thesis improves our understanding of these aspects of missing data via rigorous theoretical characterisation and extensive experimental validation of this theory. We begin by rigorously proving that more observations can hurt our ability to reconstruct signals if they are noisy, even under sample schemes designed to obviate the impact of noise, and even under robust reconstruction methods. In doing so, we enrich understanding of the  limitations of the characterisation of optimality currently used in the sampling literature. We then extend our setting to include downstream tasks, such as classification by certain classes of GNNs, rigorously characterising classification error and the relationship between reconstruction and classification error. We demonstrate the relevance of considering the downstream task in sampling and reconstruction, showing how sampling methods in the optimal setting can underperform random sampling and deriving a novel optimal sampling scheme for missing data for the task of classification.
